
\subsubsection{数据预处理}

% 仅问题 1 需要写(数据预处理)。
% 关键步骤:
%   1. 原始数据来源、格式与问题(缺失值、量纲不统一、异常值)
%   2. 数据理解与特征提取
%   3. 预处理步骤(缺失值处理 -> 标准化 -> 特征衍生),含公式
%   4. 预处理结果对比表(min/max/mean/std)
%   5. 总结

【数据来源与格式说明:原始数据来源、字段说明、规模】

\begin{enumerate}[itemindent=2em]
  \item \textbf{缺失值处理:}【方法:删除/均值填充/插值/模型预测】
  \item \textbf{异常值处理:}【方法:3$\sigma$ 准则 / IQR / 箱线图】
  \item \textbf{数据标准化:}采用 Z-score 标准化(式~\ref{eq:zscore})或 Min-Max 归一化(式~\ref{eq:minmax})。
\end{enumerate}

\begin{equation}
x'_j = \frac{x_j - \bar{x}_j}{s_j} \quad \text{(Z-score 标准化)} \label{eq:zscore}
\end{equation}

\begin{equation}
x'_j = \frac{x_j - x_{j,\min}}{x_{j,\max} - x_{j,\min}} \quad \text{(Min-Max 归一化)} \label{eq:minmax}
\end{equation}

% 处理后结果统计表(min/max/mean/std)
预处理结果统计如下表所示:

\begin{longtable}{>{\centering\arraybackslash}p{0.20\textwidth}>{\centering\arraybackslash}p{0.18\textwidth}>{\centering\arraybackslash}p{0.18\textwidth}>{\centering\arraybackslash}p{0.18\textwidth}>{\centering\arraybackslash}p{0.18\textwidth}}
  \caption{预处理后数据统计}\\
  \toprule
  变量 & 均值 & 标准差 & 最小值 & 最大值 \\
  \midrule
  \endfirsthead
  \caption{预处理后数据统计(续)}\\
  \toprule
  变量 & 均值 & 标准差 & 最小值 & 最大值 \\
  \midrule
  \endhead
  \bottomrule
  \endfoot
  【变量1】& 【值】& 【值】& 【值】& 【值】\\
  【变量2】& 【值】& 【值】& 【值】& 【值】\\
  【变量3】& 【值】& 【值】& 【值】& 【值】\\
  【变量4】& 【值】& 【值】& 【值】& 【值】\\
\end{longtable}

经过上述预处理,数据质量得到有效改善,为后续建模提供了可靠的输入。

% 图 1 【预处理前后对比图】(可选)
% \begin{figure}[ht]
%   \begin{center}
%     \includegraphics[width=0.85\textwidth]{../求解/问题1/图片/1.数据预处理对比.png}
%     \caption{预处理前后数据对比}
%     \label{fig:q1_preprocess}
%   \end{center}
% \end{figure}

综上所述,本节完成了数据预处理,经过清洗和标准化后的数据将作为后续建模的输入。
